% Бүлэг 3
\chapter{Системийн хөгжүүлэлт}
\label{Chapter3}

Энэхүү бүлэгт онлайн дуудлага худалдааны бүрэн цогц системийн хөгжүүлэлтийн талаар дэлгэрэнгүй тайлбарлана. Систем нь 3 үндсэн бүрэлдэхүүн хэсгээс бүрдэнэ: (1) Backend API систем - серверийн логик, өгөгдлийн удирдлага; (2) Веб аппликейшн - React.js ашигласан хэрэглэгчийн веб интерфэйс; (3) Мобайл аппликейшн - React Native + Expo ашигласан iOS/Android апп. Бүх бүрэлдэхүүн хэсэг нь нэг төвлөрсөн backend API-той холбогдож ажилладаг.

\section{Хөгжүүлэлтийн орчин, технологийн сонголт}

\subsection{Backend технологийн сонголт}
Backend:
\begin{itemize}
    \item Node js 
    \begin{itemize}
   \item Асинхрон, хурдан (Google V8 engine)
   \item Нэг process-д олон холболт (Event Loop)
   \item Real-time систем хөгжүүлэхэд тохиромжтой.
\end{itemize}
    \item Express js
    \begin{itemize}
   \item REST API бүтээхэд хялбар (GET, POST, PUT, DELETE)
   \item Routes (URL хаяг) удирдах нь энгийн
\end{itemize}
    \item REST API
    \begin{itemize}
   \item CRUD үйлдлүүд (Create, Read, Update, Delete) хийхэд тохиромжтой
   \item JSON формат-аар өгөгдөл дамжуулдаг (хялбар)
\end{itemize}
\end{itemize}
Database:
\begin{itemize}
    \item MongodDB  (NoSQL шинжтэй, dynamic schema-тай аппликейшнд тохиромжтой.)
\end{itemize}

 Бусад хэрэгслүүд:
 \begin{itemize}
    \item Cloudinary (зураг хадгалах үед)
    \item Github (frontend хөгжүүлэгчтэй хамтран ажиллахад)
    \item Insomnia (API туршилтанд)
\end{itemize}

\section{Системийн модуль, бүрэлдэхүүн хэсгүүдийн тайлбар}
Үнэ санал болгох модуль:
 Үнэ санал болгох модуль хэрэглэгчид бараанд үнэ санал болгох болон  дуудлага худалдааны үнийн түүх харах боломжийг олгодог бөгөөд Express router болон MongoDB-ийн bidding collection ашиглан хөгжүүлсэн.

Шинэ дуудлага худалдаа үүсгэх модуль:
 Шинэ дуудлага худалдаа үүсгэх модуль хэрэглэгчид шинэ бараа нэмэх, барааны мэдээлэл өөрчлөх, устгах боломжийг олгодог бөгөөд Express router болон MongoDB-ийн product collection ашиглан хөгжүүлсэн.

 \section{API болон өгөгдөл солилцоо}
 Манай системийн үндсэн API Endpoint-ууд:
  \begin{itemize}
    \item /api/users
    \begin{itemize}
        \item /register - Шинэ хэрэглэгч бүртгэх
        \item /login - Хэрэглэгч нэвтрэх
        \item /logout - Системээс гарах
        \item /allusers - Бүх хэрэглэгчийн жагсаалт (Админ)
        \item /userbalance - Хэрэглэгчийн дансны үлдэгдэл
        \item /send-code - Баталгаажуулах код илгээх
        \item /verify-email - Имэйл баталгаажуулах
        \item /addBalance - Данс цэнэглэх (Админ)
        \item /forgot-password - Нууц үг мартсан
        \item /verify-reset-token/:token - Token шалгах
        \item /reset-password/:token - Нууц үг шинэчлэх
         \item /google - Google OAuth нэвтрэх
         \item /google/client-id - Google Client ID авах
         \item /photo - Профайл зураг шинэчлэх

    \end{itemize}
    \item /api/product - Бараа удирдлага
    \item /api/bidding - Дуудлага худалдааны санал
    \item /api/category - Ангилал удирдлага
    \item /api/transaction - Гүйлгээний түүх
    \item /api/request - Хүсэлт удирдлага
    \item /api/likes - Таалагдсан бараа
    \item /api/watchlist - Хянах жагсаалт
    \item /api/notifications - Мэдэгдэл систем
    \item /api/reviews - Үнэлгээ, сэтгэгдэл
\end{itemize}

Бүх API endpoint-үүд JSON формат ашиглан өгөгдөл солилцдог. Authentication шаардлагатай endpoint-үүд JWT token ашиглан хэрэглэгчийг таних үйлдлийг хийнэ.

\section{Веб аппликейшны хөгжүүлэлт}
\subsection{Frontend веб технологийн сонголт}
Веб аппликейшнд дараах технологиудыг ашигласан:
\begin{itemize}
    \item \textbf{React.js}
    \begin{itemize}
   \item Component-based архитектур
   \item Virtual DOM ашиглан хурдан rendering
   \item Rich ecosystem (олон сан, plugin)
   \item Reusable компонентууд
\end{itemize}
    \item \textbf{Vite}
    \begin{itemize}
   \item Маш хурдан build tool
   \item Hot Module Replacement (HMR)
   \item Хөгжүүлэлтийн орчин хурдан ажиллана
\end{itemize}
    \item \textbf{Tailwind CSS}
    \begin{itemize}
   \item Utility-first CSS framework
   \item Responsive дизайн хийхэд хялбар
   \item Custom design system бүтээхэд тохиромжтой
\end{itemize}
    \item \textbf{React Router}
    \begin{itemize}
   \item Client-side routing
   \item Single Page Application (SPA)
   \item Хуудас шилжих хурд өндөр
\end{itemize}
    \item \textbf{Socket.io-client}
    \begin{itemize}
   \item Бодит цагийн мэдээлэл авах
   \item Дуудлага худалдааны статус шууд харагдана
\end{itemize}
\end{itemize}

\subsection{Веб аппликейшны үндсэн функцууд}
\begin{itemize}
    \item Админы удирдлагын dashboard
    \item Хэрэглэгчийн бүртгэл, нэвтрэх (Google OAuth, email)
    \item Бараа хайх, шүүлтүүр тавих
    \item Дуудлага худалдаанд оролцох (бодит цагийн мэдээлэл)
    \item Барааны дэлгэрэнгүй мэдээлэл үзэх
    \item Профайл удирдах, гүйлгээний түүх
    \item Мэдэгдлийн систем
    \item Responsive дизайн (бүх төхөөрөмж дээр)
\end{itemize}

 \section{Мобайл аппликейшны хөгжүүлэлт}
\subsection{Мобайл технологийн сонголт}
Мобайл аппликейшнд дараах технологиудыг ашигласан:
\begin{itemize}
    \item \textbf{React Native}
    \begin{itemize}
   \item Cross-platform (iOS, Android) хөгжүүлэлт нэг code base-ээр
   \item JavaScript/TypeScript ашиглан хөгжүүлдэг
   \item Native компонентуудтай ажиллах боломжтой
   \item React.js-тэй адилхан архитектур (Component-based)
\end{itemize}
    \item \textbf{Expo}
    \begin{itemize}
   \item React Native аппликейшнийг хөгжүүлэх, тест хийх хялбар орчин
   \item Hot reload - Код солих бүрт шууд харагдах
   \item Олон built-in library (Camera, Location, Notification гэх мэт)
   \item iOS болон Android дээр тест хийхэд хялбар
\end{itemize}
    \item \textbf{Expo Router}
    \begin{itemize}
   \item File-based routing систем
   \item TypeScript дэмжлэгтэй
   \item Deep linking автоматаар дэмждэг
\end{itemize}
\end{itemize}

\subsection{Мобайл аппликейшны архитектур}
Мобайл апп нь дараах үндсэн бүтэцтэй:
\begin{itemize}
    \item \textbf{Tab Navigation}: Үндсэн навигаци (Home, Search, Selling, Profile)
    \item \textbf{Stack Navigation}: Дэлгэрэнгүй хуудсууд (Product details, Category view)
    \item \textbf{Context API}: Global state management (Auth, Language)
    \item \textbf{REST API Integration}: Backend-тэй холбогдох
    \item \textbf{Real-time Updates}: Socket.io ашиглан бодит цагийн мэдээлэл
\end{itemize}

\subsection{Мобайл аппликейшны үндсэн функцууд}
\subsubsection{Хэрэглэгчийн баталгаажуулалт}
\begin{itemize}
    \item Google OAuth нэвтрэх
    \item Утасны дугаараар баталгаажуулалт (Firebase Phone Auth)
    \item JWT token ашиглан session удирдлага
    \item Автомат нэвтрэлт (AsyncStorage)
\end{itemize}

\subsubsection{Бараа удирдах}
\begin{itemize}
    \item Бараа нэмэх (20 хүртэл зураг)
    \item AI категори санал болгох систем
    \item Rich text editor ашиглан дэлгэрэнгүй тайлбар
    \item Автомашины тусгай мэдээлэл (VIN, Model, Year гэх мэт)
    \item Дуудлага худалдаа эхлэх огноо тохируулах
\end{itemize}

\subsubsection{Дуудлага худалдаа}
\begin{itemize}
    \item Бодит цагийн үнийн шинэчлэлт
    \item Countdown таймер
    \item Үнийн түүхийн график
    \item Өрсөлдөгчдийн мэдэгдэл
    \item Watchlist функц
\end{itemize}

\subsection{Мобайл дизайн загвар}
Мобайл апп нь Mongolian marketplace-д зориулагдсан учир:
\begin{itemize}
    \item Монгол, Англи хэлний дэмжлэг
    \item Төгрөг валютын харуулалт
    \item Японы аукшин сайтуудын (Mercari, Yahoo Auctions) дизайн загвар
    \item Responsive дизайн (Phone, Tablet)
    \item Dark mode дэмжлэг
\end{itemize}

 \section{Аюулгүй байдал ба нэвтрэлт}
Манай системд хэрэглэгч нэвтрэхдээ Админ болон хэрэглэгч гэж 2 түвшинд хуваагддаг.
\begin{itemize}
\item \textbf{Энгийн хэрэглэгч}: Дуудлагад оролцох, бараа байршуулах эрхтэй
\item \textbf{Админ хэрэглэгч}: Системийн бүх үйл ажиллагааг удирдах эрхтэй
\end{itemize}

\subsection{Ашиглах технологи}
\begin{itemize}
\item \textbf{Express.js middleware}: Хандалтыг шалгах
\item \textbf{JWT токен}: Хэрэглэгчийн session удирдах
\item \textbf{MongoDB}: Хэрэглэгчийн мэдээллийг хадгалах
\end{itemize}
\subsection{Ажиллах процесс}
\begin{enumerate}
\item Хэрэглэгч нэвтрэхдээ токен үүсгэгдэнэ
\item Хүсэлт бүрд токеныг шалгана
\item Токен хүчингүй бол хандалтыг хязгаарлана
\end{enumerate}

\subsection{Алдааны Тохиолдолд}

\begin{itemize}
\item Токен байхгүй: \texttt{401 Unauthorized}
\item Админ эрхгүй: \texttt{403 Forbidden}
\item Буруу токен: \texttt{401 Unauthorized}
\end{itemize}
Манай систем нь өгөгдлүүдээ mongoDB дээр хадгалдаг  ба зураг файлыг cloudinary-д байршуулдаг.Нууц үг гэх мэт эмзэг өгөгдлийг bcrypt ашиглаж шифрлэлт хийж хадгалдаг .
 

% \section{Бүлгийн дүгнэлт}
% Зохиомжийн хэсэгт өгөгдлийн ерөнхий схем болон өмнөх шинжилгээний класс диаграм дээр үндэслэн зохиомжийн класс диаграм, дарааллын диаграмуудыг гаргасан. Хамгийн гол төлөвийн диаграмыг бас давхар гаргаж өгсөн байгаа. Бусад ижил төстэй томоохон сайтуудаас жишээ авч вебийн туршилтын загварыг гаргасан байгаа.
